\documentclass[11pt]{article}
\usepackage[margin=1in]{geometry}
\usepackage{amsmath,amssymb,amsthm,mathtools}
\usepackage{graphicx}
\usepackage[hidelinks]{hyperref}
\usepackage{microtype}
\usepackage{enumitem}

\newtheorem{theorem}{Theorem}
\newtheorem{corollary}[theorem]{Corollary}
\newtheorem{proposition}[theorem]{Proposition}
\newtheorem{remark}[theorem]{Remark}
\newcommand{\C}{\mathbb C}
\newcommand{\M}{\mathsf M}
\newcommand{\T}{\mathbb T}
\newcommand{\rank}{\operatorname{rank}}
\newcommand{\Ran}{\operatorname{Ran}}
\newcommand{\Span}{\operatorname{span}}

\title{Normal Targeting at Corank One:\\A Complete One-Column Criterion}
\author{Automated mathematical discovery packet\\Run \texttt{fa\_banach\_001}, lane 8}
\date{27 August 2026}

\begin{document}
\maketitle

\begin{center}
\fbox{\parbox{0.90\linewidth}{\textbf{Status.} Likely valid partial result,
subject to human review.  The general normal targeting problem remains open;
this packet completely treats the subcase $\rank X=m-1$ over $\C$.}}
\end{center}

\begin{abstract}
Given $X,Y\in\M_{m\times n}(\C)$, the normal targeting problem asks when a
normal $A\in\M_m(\C)$ satisfies $AX=Y$.  We solve the maximal-rank singular
case $\rank X=m-1$.  A compact singular value decomposition reduces the
problem to completing one prescribed block column
$\left[\begin{smallmatrix}B\\c^*\end{smallmatrix}\right]$ by one free column.
Normality is equivalent to a positive rank-one defect identity and a single
vector compatibility equation involving only one phase and one complex
scalar.  When $B$ is normal, the latter condition says exactly that the
eigenvalues of $B$ visible from $c$ lie on an affine line.  This gives an
explicit $4\times4$ obstruction showing that the positive-semidefinite defect
condition alone is not sufficient even when only one column is missing.
\end{abstract}

\section{Source question and scope}

Bierly, Garcia, and Horn introduce the linear targeting problem and state in
Section 10 that the normal case remains open \cite{BGH}.  The relevant source
passage (p.~13) is reproduced below.

\begin{center}
\includegraphics[width=0.96\linewidth]{figures/open_problem_crop.png}
\end{center}

The question is:
\begin{quote}
Characterize those $X,Y\in\M_{m\times n}(\C)$ for which there is a normal
$A\in\M_m(\C)$ satisfying $AX=Y$.
\end{quote}
The source solves several structured variants, including normal matrices with
two prescribed spectral values, but not the unrestricted problem.  We treat
the first rank below the invertible case: $\rank X=m-1$.  Since the source
assumes $m\ge n$, this forces $n\in\{m-1,m\}$.

\section{Proof intuition}

The columns of $X$ determine the action of $A$ on the $(m-1)$-dimensional
space $\Ran X$.  After choosing an orthonormal basis adapted to this space,
every targeting matrix therefore has all but its last column fixed.  Write it
as
\[
 \widehat A=\begin{bmatrix}B&d\\ c^*&\eta\end{bmatrix}.
\]
Comparing the upper-left and upper-right blocks of
$\widehat A^*\widehat A$ and $\widehat A\widehat A^*$ gives the whole answer.
The upper-left comparison says that the computable defect
$B^*B-BB^*+cc^*$ must be $dd^*$, hence positive semidefinite of rank at most
one.  If it has rank one, it determines $d$ up to a phase.  The upper-right
comparison is then one explicit compatibility equation for that phase and
$\eta$.  The bottom-right comparison follows automatically by taking traces.

When $B$ is normal, the defect forces $d=\omega c$.  In a spectral basis for
$B$, the compatibility equation says that
$\lambda-\omega\overline\lambda$ is constant on every eigenvalue seen by
$c$.  Geometrically this means precisely that those eigenvalues lie on one
line in the complex plane.

\section{The corank-one criterion}

Let $r=m-1$.  Take a compact singular value decomposition
\[
 X=V_1\Sigma W_1^*,
\]
where $V_1\in\M_{m\times r}(\C)$ and
$W_1\in\M_{n\times r}(\C)$ have orthonormal columns and
$\Sigma\in\M_r(\C)$ is positive diagonal.  Complete $V_1$ to a unitary
$V=[\,V_1\ v\,]$.  Assuming $\ker X\subseteq\ker Y$, set
\[
 Z:=YW_1\Sigma^{-1},\qquad
 V^*Z=\begin{bmatrix}B\\c^*\end{bmatrix},
 \quad B\in\M_r(\C),\ c\in\C^r,
\]
and define the Hermitian defect
\[
 H:=B^*B-BB^*+cc^*.
\]

\begin{theorem}[Complete criterion when $\rank X=m-1$]\label{thm:main}
Let $X,Y\in\M_{m\times n}(\C)$, $m\ge n$, and $\rank X=m-1$.
There is a normal $A\in\M_m(\C)$ satisfying $AX=Y$ if and only if:
\begin{enumerate}[label=\textup{(\roman*)}]
\item $\ker X\subseteq\ker Y$;
\item $H$ is positive semidefinite and $\rank H\le1$; and
\item if $\rank H=1$ and $h\in\C^r$ is any vector with $H=hh^*$,
there are $\omega\in\T$ and $\eta\in\C$ such that
\begin{equation}\label{eq:compat}
 B^*(\omega h)+c\eta=Bc+(\omega h)\overline\eta.
\end{equation}
\end{enumerate}
If $H=0$, condition \textup{(iii)} is vacuous.  In the rank-one case, every
normal targeting matrix is obtained from a solution $(\omega,\eta)$ of
\eqref{eq:compat} by
\begin{equation}\label{eq:construction}
 A=V\begin{bmatrix}B&\omega h\\c^*&\eta\end{bmatrix}V^*.
\end{equation}
When $H=0$, every $\eta\in\C$ gives a solution by using the same formula with
the upper-right vector equal to zero.
\end{theorem}

\begin{proof}
The inclusion $\ker X\subseteq\ker Y$ is necessary for any targeting matrix.
It is equivalent to $\Ran Y^*\subseteq\Ran X^*=\Ran W_1$, and hence
\begin{equation}\label{eq:Ysupport}
 Y=YW_1W_1^*.
\end{equation}
Moreover, $AX=Y$ is equivalent to $AV_1=Z$.  Indeed, $AX=Y$ gives this after
right multiplication by $W_1\Sigma^{-1}$, while the converse follows from
\eqref{eq:Ysupport}:
\[
 AX=AV_1\Sigma W_1^*=Z\Sigma W_1^*=YW_1W_1^*=Y.
\]
Consequently every targeting matrix, and only a targeting matrix, has the
form
\[
 A=V\widehat A V^*,\qquad
 \widehat A=\begin{bmatrix}B&d\\c^*&\eta\end{bmatrix}
\]
for some $d\in\C^r$ and $\eta\in\C$.

Direct multiplication gives
\[
 \widehat A^*\widehat A=
 \begin{bmatrix}
 B^*B+cc^*&B^*d+c\eta\\
 d^*B+\overline\eta c^*&d^*d+|\eta|^2
 \end{bmatrix}
\]
and
\[
 \widehat A\widehat A^*=
 \begin{bmatrix}
 BB^*+dd^*&Bc+d\overline\eta\\
 c^*B^*+\eta d^*&c^*c+|\eta|^2
 \end{bmatrix}.
\]
Thus the upper-left blocks agree exactly when
\begin{equation}\label{eq:defect}
 H=dd^*,
\end{equation}
and the upper-right blocks agree exactly when
\begin{equation}\label{eq:cross}
 B^*d+c\eta=Bc+d\overline\eta.
\end{equation}
Taking traces in \eqref{eq:defect}, and using
$\operatorname{tr}(B^*B-BB^*)=0$, yields $d^*d=c^*c$; hence the bottom-right
blocks then agree automatically.  The lower-left equality is the adjoint of
\eqref{eq:cross}.  Therefore \eqref{eq:defect} and \eqref{eq:cross} are
necessary and sufficient for normality.

Equation \eqref{eq:defect} is solvable precisely when $H\ge0$ and
$\rank H\le1$.  If $H=hh^*$ has rank one, all of its solutions are
$d=\omega h$ with $\omega\in\T$, and \eqref{eq:cross} becomes
\eqref{eq:compat}.  If $H=0$, then $d=0$; taking the trace also gives $c=0$,
and $H=0$ reduces to $B^*B=BB^*$.  Hence \eqref{eq:cross} is automatic for
every $\eta$.  This proves the criterion and construction.  Unitary
conjugation by $V$ preserves normality.
\end{proof}

\begin{remark}[A three-real-variable test]
Once $H$ has passed the rank-one test, the original $m^2$ normality equations
have been reduced to \eqref{eq:compat}: one unit-modulus phase $\omega$ and
one complex scalar $\eta$.  For fixed $\omega$, this is a real-linear system
in $\Re\eta$ and $\Im\eta$.  Thus the criterion is directly checkable and
constructive without optimizing over a full $m\times m$ matrix.
\end{remark}

\section{Closed form for a normal compression}

Suppose now that $B$ is normal.  Let
\[
 B=U\operatorname{diag}(\lambda_1,\ldots,\lambda_r)U^*,
 \qquad \widehat c=U^*c,
\]
and call
\[
 \Lambda_c:=\{\lambda_j:\widehat c_j\ne0\}
\]
the \emph{active spectrum} of $(B,c)$.  Equivalently, an eigenvalue is active
when the spectral projection of $B$ at that eigenvalue does not annihilate
$c$; this definition is independent of the chosen basis within repeated
eigenspaces.

\begin{corollary}[Affine-line criterion]\label{cor:line}
If $B$ is normal, the prescribed block column
$\left[\begin{smallmatrix}B\\c^*\end{smallmatrix}\right]$ has a normal
one-column completion if and only if either $c=0$ or $\Lambda_c$ is contained
in an affine line in $\C$.
\end{corollary}

\begin{proof}
For normal $B$, $H=cc^*$.  When $c\ne0$, Theorem~\ref{thm:main} permits
$h=c$, so $d=\omega c$ for some $\omega\in\T$.  In the spectral basis of
$B$, equation \eqref{eq:compat} says, for every active index $j$,
\begin{equation}\label{eq:spectral}
 \lambda_j-\omega\overline{\lambda_j}
 =\eta-\omega\overline\eta.
\end{equation}
Write $\omega=e^{2i\theta}$.  Multiplication by $e^{-i\theta}$ transforms
the left side of \eqref{eq:spectral} into
\[
 e^{-i\theta}\lambda_j-e^{i\theta}\overline{\lambda_j}
 =2i\,\operatorname{Im}(e^{-i\theta}\lambda_j).
\]
It is independent of $j$ precisely when all active eigenvalues have the same
orthogonal coordinate relative to direction $e^{i\theta}$, that is, precisely
when they lie on an affine line parallel to $e^{i\theta}$.  If such a line
exists, choose $\eta$ on the same line; then \eqref{eq:spectral} holds and the
construction in Theorem~\ref{thm:main} gives a normal completion.  The case
$c=0$ is immediate.
\end{proof}

\section{A sharp four-dimensional obstruction}

\begin{proposition}\label{prop:counterexample}
Let
\[
 X=\begin{bmatrix}1&0&0\\0&1&0\\0&0&1\\0&0&0\end{bmatrix},\qquad
 Y=\begin{bmatrix}0&0&0\\0&1&0\\0&0&i\\1&1&1\end{bmatrix}.
\]
Then $\rank X=3$, $\ker X\subseteq\ker Y$, and the source paper's necessary
defect matrix $H(B,C)$ is positive semidefinite of rank one, but there is no
normal $A\in\M_4(\C)$ with $AX=Y$.
\end{proposition}

\begin{proof}
Here $V=I_4$, $\Sigma=I_3$,
\[
 B=\operatorname{diag}(0,1,i),\qquad c=(1,1,1)^T.
\]
Since $B$ is normal,
\[
 H=B^*B-BB^*+cc^*=cc^*\ge0,
 \qquad \rank H=1.
\]
However, the active spectrum is $\{0,1,i\}$, which is not collinear, so
Corollary~\ref{cor:line} rules out a normal completion.

For a direct algebraic check, normality would force $d=\omega c$ with
$\omega\in\T$ and
\[
 \lambda-\omega\overline\lambda
 =\eta-\omega\overline\eta
 \quad(\lambda\in\{0,1,i\}).
\]
The equation for $\lambda=0$ makes the common value zero.  The equation for
$\lambda=1$ then forces $\omega=1$, while the equation for $\lambda=i$
forces $\omega=-1$, a contradiction.
\end{proof}

The source already observes that its positive-semidefinite defect condition
is not sufficient for general square border blocks of size at least three.
Proposition~\ref{prop:counterexample} sharpens the phenomenon to a single
missing column and gives the complete geometric reason in the normal-$B$
case.  Within this family, three active noncollinear eigenvalues are necessary,
so dimension four is minimal.

\section{Verification notes}

\begin{itemize}[leftmargin=*]
\item The proof uses only compact SVD, direct block multiplication, and the
spectral theorem for the normal-compression corollary.
\item The potentially delicate $n=m$ case is covered by
$\ker X\subseteq\ker Y$: it gives $Y=YW_1W_1^*$ even though $W_1$ is not
square.
\item The trace of the upper-left defect identity gives
$\|d\|=\|c\|$, so no independent bottom-right condition is missing.
\item The deterministic script \texttt{code/verify\_codim\_one.py} checks
the two block residuals on a random unitarily diagonalized normal matrix,
constructs a non-Hermitian positive example with collinear active spectrum,
and verifies the rank-one defect and incompatible phases in
Proposition~\ref{prop:counterexample}.
\end{itemize}

\section{Limitations and novelty audit}

This is a partial result: it does not characterize normal targeting when
$\operatorname{codim}\Ran X\ge2$.  In the nonnormal-$B$ case, the theorem
leaves a one-phase compatibility test rather than eliminating that phase by a
closed formula.

A bounded search on 27 August 2026 used the exact arXiv id and title, the
phrases ``normal targeting'', ``normal completion one missing column'', and
``normal defect one'', and the authors' names.  It found the source paper and
the established normal-defect-one literature, especially
Kaliuzhnyi-Verbovetskyi, Spitkovsky, and Woerdeman \cite{KSW}.  That work
classifies completions of a prescribed north-west block when both border
vectors may vary.  Our targeting reduction instead prescribes the lower
border $c^*$ and leaves only the upper border and corner free.  No source found
in the bounded search stated the rank-$m-1$ targeting criterion, the active
spectrum affine-line criterion, or the explicit $4\times4$ obstruction above.
Absolute novelty is not claimed; the result is close in spirit to normal
completion theory and may be derivable from a broader formulation not found
in the search.

\section{Human-review recommendation}

Recommend expert review as a likely valid partial solution.  The two points
most worth checking are (1) the row-space implication
$Y=YW_1W_1^*$ under $\ker X\subseteq\ker Y$, and (2) the conjugations in
equation \eqref{eq:compat}.  If those pass, the theorem gives an exact and
constructive solution of the maximal-rank singular subcase and the
four-dimensional example is a rigorous obstruction, not a numerical claim.

\enlargethispage{3\baselineskip}
\small
\begin{thebibliography}{9}
\bibitem{BGH}
K. Bierly, S. R. Garcia, and R. A. Horn,
\emph{The linear targeting problem}, arXiv:2408.10036, 2024 (revised 2025).

\bibitem{KSW}
D. S. Kaliuzhnyi-Verbovetskyi, I. M. Spitkovsky, and H. J. Woerdeman,
\emph{Matrices with normal defect one}, Operators and Matrices
\textbf{3} (2009), 427--450; arXiv:0903.0090.
\end{thebibliography}

\end{document}
